\begin{theorem}[No Object Is Its Own Successor]
\label{thm:no-object-is-its-own-successor}
Let $(P,S,1)$ be a Peano system. For every $n \in P$,
\[
S(n) \neq n.
\]
\hyperref[prf:no-object-is-its-own-successor]{Go to proof.}
\end{theorem}
\end{theorembox}

\begin{formalizationrecord}{Lean 4 Verification Record}
\FormalizationSource{Landau, \emph{Foundations of Analysis}, Section 2, Theorem 2}
\Formalizes{thm:no-object-is-its-own-successor}
\textbf{Lean module:} \texttt{LRA.VolumeII.PeanoSystems.PeanoSystem}.\par
\LeanDeclaration{NoObjectIsItsOwnSuccessor}
\textbf{Status:} pending against \texttt{lra-lean}.\par

\begin{leancode}
theorem NoObjectIsItsOwnSuccessor
    (ps : PeanoSystem) :
    forall element : ps.carrier,
      ps.successor element ≠ element := by
  sorry
\end{leancode}
\end{formalizationrecord}

\begin{remark*}[Standard quantified statement]
For a Peano system $(P,S,1)$,
\[
\forall n \in P\;(S(n) \neq n).
\]
\end{remark*}

\begin{remark*}[Interpretation]
Successor always moves forward in the counting structure. No element can be
recovered by applying successor to itself.
\end{remark*}

\begin{remark*}[Historical note]
This theorem is a standard consequence of the Peano axioms in the classical
development of the natural numbers \citep{PeanoArithmeticesPrincipia1889}.
Feferman includes the corresponding result as an exercise in Chapter 3 of
\citet{FefermanNumberSystems1964}.
\end{remark*}

\begin{workedexample}[Why the base case is forced]
\LRAWorkedExampleFor{thm:no-object-is-its-own-successor}
\LRAWorkedExampleUses{ax:peano-base-not-successor}
\LRAWorkedExampleTags{peano systems, induction, base case}
For the first element, the theorem says
\[
S(1)\neq 1.
\]
This is exactly the axiom that the distinguished element is not a successor:
there is no element $x\in P$ with $S(x)=1$. In particular, taking $x=1$
would be impossible, so $S(1)\neq 1$.
\end{workedexample}

\begin{workedexample}[Why the induction step uses injectivity]
\LRAWorkedExampleFor{thm:no-object-is-its-own-successor}
\LRAWorkedExampleUses{ax:peano-successor-injective, thm:induction-principle-for-peano-system}
\LRAWorkedExampleTags{peano systems, induction, injectivity}
Suppose the induction hypothesis gives $S(n)\neq n$. To prove the next case,
we must show
\[
S(S(n))\neq S(n).
\]
If equality held, injectivity of $S$ would let us cancel the outer successor
and conclude $S(n)=n$, contradicting the induction hypothesis. Thus the
property propagates from $n$ to $S(n)$.
\end{workedexample}

\begin{dependencies}
\begin{itemize}
  \item \hyperref[ax:peano-base-not-successor]{Distinguished Element Is Not a Successor}
  \item \hyperref[ax:peano-successor-injective]{Injectivity of Successor}
  \item \hyperref[thm:induction-principle-for-peano-system]{Induction Principle for a Peano System}
\end{itemize}
\end{dependencies}
\end{topicbox}
